\documentclass[../main.tex]{subfiles}
\begin{document}
\chapter{Constraint Velocities and Accelerations}
\label{chap:constraint-kinematics}

\section{State of the Art}
Constraint kinematics follow by differentiating holonomic position equations. Evaluation of the acceleration bias and Jacobian derivatives is important for simulation and reaction-force calculations \cite{shabana2013dynamics}. This focused review uses the standard analytical derivative hierarchy for the planar pin model; it does not compare automatic-differentiation or recursive implementations. Bibliographic scope checked 2026-10-07.

\section{Position constraints imply velocity constraints}

Begin with the assembled holonomic equations
\begin{equation}
	\vect{\Phi}(\vect{q})=\vect{0}.
	\label{eq:holonomic-position}
\end{equation}
Define the constraint Jacobian $J(\vect{q})=\partial\vect{\Phi}/\partial\vect{q}$. Differentiating~\eqref{eq:holonomic-position} with respect to time gives
\begin{equation}
	\dot{\vect{\Phi}}=J(\vect{q})\dot{\vect{q}}=\vect{0}.
	\label{eq:velocity-constraint}
\end{equation}
The admissible generalized velocities therefore lie in the null space of $J$. At a regular configuration, that null space has dimension $n_q-\operatorname{rank}(J)$, equal to the local mobility.

\section{Acceleration constraints and curvature terms}

Differentiate~\eqref{eq:velocity-constraint} again. Because $J$ depends on configuration, its time derivative contributes:
\begin{equation}
	\ddot{\vect{\Phi}}
	=J(\vect{q})\ddot{\vect{q}}
	 +\dot J(\vect{q},\dot{\vect{q}})\dot{\vect{q}}
	=\vect{0}.
	\label{eq:acceleration-constraint}
\end{equation}
The vector $\dot J\dot{\vect{q}}$ is sometimes called the acceleration bias or curvature term. It is generally nonzero even though the position constraints are homogeneous and time-independent.

For a point fixed in a rotating planar body, write its body-frame offset as $\vect{s}$ and let $\vect{\rho}=A(\theta)\vect{s}$ be the rotated offset. Since $\dot{\vect{\rho}}=\dot\theta S\vect{\rho}$, its world acceleration is
\begin{equation}
	\ddot{\vect{r}}
	=\ddot{\vect{p}}+\ddot\theta S\vect{\rho}-\dot\theta^2\vect{\rho}.
	\label{eq:rotating-point-acceleration}
\end{equation}
The last term is centripetal acceleration. For a pin residual $\vect{r}_a-\vect{r}_b$, take the signed difference of the two attached-point accelerations. The terms multiplying translational and angular accelerations form $J\ddot{\vect{q}}$; the remaining squared-angular-speed terms form $\dot J\dot{\vect{q}}$. Fixed-pose constraints are linear in the coordinates and contribute no curvature term.
If $\vect\rho_a=A(\theta_a)\vect s_a$ and $\vect\rho_b=A(\theta_b)\vect s_b$ are the world attachment offsets, the two-component pin bias is
\begin{equation}
\vect\gamma_{ab}=-\dot\theta_a^2\vect\rho_a+\dot\theta_b^2\vect\rho_b.
\label{eq:pin-acceleration-bias}
\end{equation}
The signs follow from subtracting the centripetal terms $-\dot\theta_i^2\vect\rho_i$ in the two point accelerations according to the first-body-minus-second-body residual convention.

\section{Five-link constraint derivative hierarchy}

For the five-link model, the complete hierarchy is
\begin{equation}
	\begin{aligned}
		\text{position:}\quad &\vect{\Phi}(\vect{q})=\vect{0},\\
		\text{velocity:}\quad &J(\vect{q})\dot{\vect{q}}=\vect{0},\\
		\text{acceleration:}\quad &J(\vect{q})\ddot{\vect{q}}
			+\vect{\gamma}(\vect{q},\dot{\vect{q}})=\vect{0},
	\end{aligned}
	\label{eq:constraint-hierarchy}
\end{equation}
where $\vect{\gamma}=\dot J\dot{\vect{q}}$. The symbolic notebook derives this term by differentiating the velocity-level equations; the numerical model independently computes the same centripetal contributions. Their agreement is checked at the projected initial state before integration.
The notebook computes it in two independent ways. Symbolically, form $\dot{\vect\Phi}=J\dot{\vect q}$, differentiate with respect to $\vect q$ and multiply by $\dot{\vect q}$; this is the part of the total derivative that excludes $J\ddot{\vect q}$. Numerically, evaluate~\eqref{eq:pin-acceleration-bias} for each pin and append zeros for the fixed-pose rows. Compare these vectors at the same consistent state before integration. At the notebook's projected initial state, the maximum component difference is $4.337\times10^{-19}$.

For the symbolic notebook's midpoint links, each pin residual joins the right endpoint of link $i$ to the left endpoint of link $j$. Write $c_i=\cos\theta_i$, $s_i=\sin\theta_i$, and $\omega_i=\dot\theta_i$. The 13 ideal acceleration residuals, ordered as the two rows for each pin followed by the grounded pose, are
\begin{equation}
\begin{aligned}
r_{12,x}={}&\ddot x_1-\ddot x_2
-\frac{L_1}{2}(\ddot\theta_1s_1+\omega_1^2c_1)
-\frac{L_2}{2}(\ddot\theta_2s_2+\omega_2^2c_2)=0,\\
r_{12,y}={}&\ddot y_1-\ddot y_2
+\frac{L_1}{2}(\ddot\theta_1c_1-\omega_1^2s_1)
+\frac{L_2}{2}(\ddot\theta_2c_2-\omega_2^2s_2)=0,\\[2pt]
r_{23,x}={}&\ddot x_2-\ddot x_3
-\frac{L_2}{2}(\ddot\theta_2s_2+\omega_2^2c_2)
-\frac{L_3}{2}(\ddot\theta_3s_3+\omega_3^2c_3)=0,\\
r_{23,y}={}&\ddot y_2-\ddot y_3
+\frac{L_2}{2}(\ddot\theta_2c_2-\omega_2^2s_2)
+\frac{L_3}{2}(\ddot\theta_3c_3-\omega_3^2s_3)=0,\\[2pt]
r_{34,x}={}&\ddot x_3-\ddot x_4
-\frac{L_3}{2}(\ddot\theta_3s_3+\omega_3^2c_3)
-\frac{L_4}{2}(\ddot\theta_4s_4+\omega_4^2c_4)=0,\\
r_{34,y}={}&\ddot y_3-\ddot y_4
+\frac{L_3}{2}(\ddot\theta_3c_3-\omega_3^2s_3)
+\frac{L_4}{2}(\ddot\theta_4c_4-\omega_4^2s_4)=0,\\[2pt]
r_{45,x}={}&\ddot x_4-\ddot x_5
-\frac{L_4}{2}(\ddot\theta_4s_4+\omega_4^2c_4)
-\frac{L_5}{2}(\ddot\theta_5s_5+\omega_5^2c_5)=0,\\
r_{45,y}={}&\ddot y_4-\ddot y_5
+\frac{L_4}{2}(\ddot\theta_4c_4-\omega_4^2s_4)
+\frac{L_5}{2}(\ddot\theta_5c_5-\omega_5^2s_5)=0,\\[2pt]
r_{51,x}={}&\ddot x_5-\ddot x_1
-\frac{L_5}{2}(\ddot\theta_5s_5+\omega_5^2c_5)
-\frac{L_1}{2}(\ddot\theta_1s_1+\omega_1^2c_1)=0,\\
r_{51,y}={}&\ddot y_5-\ddot y_1
+\frac{L_5}{2}(\ddot\theta_5c_5-\omega_5^2s_5)
+\frac{L_1}{2}(\ddot\theta_1c_1-\omega_1^2s_1)=0,\\[2pt]
r_{g,x}={}&\ddot x_4=0,\qquad
r_{g,y}=\ddot y_4=0,\qquad
r_{g,\theta}=\ddot\theta_4=0.
\end{aligned}
\label{eq:five-link-expanded-acceleration-residuals}
\end{equation}
Collecting these entries gives exactly $J\ddot{\vect q}+\dot J\dot{\vect q}=\vect0$. The terms proportional to $\ddot\theta_i$ come from tangential acceleration of the endpoint offsets; the $\omega_i^2$ terms are centripetal contributions. In the notebook's stabilized solve, Baumgarte feedback is formed over the full 13-entry residual, including the three grounded-pose rows. Those final three correction terms are zero while link~4 remains exactly at its reference pose with zero velocity; they are not omitted from the formula.

Equation~\eqref{eq:acceleration-constraint} supplies the kinematic equations needed to determine accelerations compatible with loop closure. In the next chapter, these equations are coupled to force balance, with Lagrange multipliers representing ideal joint reactions.

\section{Summary}

Differentiating holonomic constraints produces velocity compatibility and then acceleration compatibility. The acceleration equation contains both the unknown generalized accelerations and the configuration/velocity-dependent bias term. Omitting that term would fail to enforce the actual acceleration of rotating pin locations.

\ifSubfilesClassLoaded{\backmatter\bookbibliography}{}
\end{document}